We are going to learn to program a computer by making it do
drawings. To start simply,
we use squared paper with small squares.
The {\bf drawing rule} \index{Drawing rule}
to be respected is that {\em our drawings may only contain straight
lines and quarter circles that must join
two ends
of the same horizontal side,
of the same vertical side or
of the same diagonal
of a square of the grid}. This drawing rule therefore
gives us the possibility of drawing:
\begin{list}{$\scriptscriptstyle\bullet$}{
}
\item the horizontal segments
\begin{picture}(12,10)
\put(0,0){\line(1,0){10}}
\end{picture} \/ corresponding to the horizontal sides of a square of the
grid;
\item the vertical segments
\begin{picture}(6,10)
\put(2,0){\line(0,1){10}}
\end{picture} corresponding to the vertical sides of a square of the
grid;
\item the oblique segments
\begin{picture}(10,10)
\put(0,0){\line(1,1){10}}
\end{picture}
and
\begin{picture}(10,10)
\put(0,10){\line(1,-1){10}}
\end{picture}\/ corresponding to the diagonals of a square of the
grid;
\item the quarter circles \/
\begin{picture}(10,10)
\put(0,0){\makebox(10,10)[bl]{$\displaystyle \smile$}}
\end{picture} \/ and \/
\begin{picture}(10,10)
\put(0,0){\makebox(10,10)[bl]{$\displaystyle \frown$}}
\end{picture} \/ joining the ends of the same horizon\-tal side of a square of the
grid;
\item the quarter circles
\begin{picture}(5,10)
\put(0,0){\makebox(5,10)[bl]{\footnotesize (}}
\end{picture} and
\begin{picture}(5,10)
\put(0,0){\makebox(5,10)[bl]{\footnotesize )}}
\end{picture} joining the ends of the same vertical side of a square of the
grid;
\item the quarter circles
\begin{picture}(10,10)
\put(10,0){\oval(20,20)[tl]}
\end{picture}\ ,
\begin{picture}(10,10)
\put(0,10){\oval(20,20)[br]}
\end{picture}\ ,
\begin{picture}(10,10)
\put(0,0){\oval(20,20)[tr]}
\end{picture} \/
and \/
\begin{picture}(10,10)
\put(10,10){\oval(20,20)[bl]}
\end{picture} \/ joining the ends of the same diagonal of a square
of the grid.
\end{list}

\noindent The small drawings of figure \ref{fig:pet-des-geo} (page
\pageref{fig:pet-des-geo}) are simple examples drawn while
respecting our drawing rule on
squared paper.
\begin{figure}[htbp]
\includevspacePostScript[2.4cm]{petits-dessins}{0.95}{0.95}
\caption{Some simple geometric drawings}
\label{fig:pet-des-geo}
\end{figure}
Although this rule does not leave complete
freedom to our imagination,
we can
make figurative drawings such as those of the stegosaurus or of the
fortified castle on page \pageref{fig:stegosaure}.
\begin{figure}[htbp]
 \includevspacePostScript[6.1cm]{stegosaure}{0.95}{0.95}
 \caption{Stegosaurus}
 \label{fig:stegosaure}
\end{figure}
\begin{figure}[htbp]
 \includevspacePostScript[7.4cm]{chateaufort}{0.95}{0.95}
 \caption{Fortified castle}\label{fig:chateaufort}
\end{figure}

\pagebreak[4]
\needspace{6\baselineskip}\noindent \rule{137mm}{0.1mm}\par\nopagebreak %---------------------------------------------
\begin{exercices}{\ \\

\vspace{-\baselineskip}
\tiret Complete the drawings of figure \ref{fig:completer-les-dessins}
below, using only
horizontal
\begin{picture}(12,10)
\put(0,0){\line(1,0){10}}
\end{picture},
vertical
\begin{picture}(6,10)
\put(2,0){\line(0,1){10}}
\end{picture}
or oblique
\begin{picture}(10,10)
\put(0,0){\line(1,1){10}}
\end{picture}
and
\begin{picture}(10,10)
\put(0,10){\line(1,-1){10}}
\end{picture}
lines, which must join two vertices of the grid with
successive numbers, the vertices
being shown on this figure by small dots.

\tiret Make geometric drawings on squared paper
with small squares, respecting the drawing rule.
Here are a few ideas (these drawings will later be made on the com\-put\-er):
a hopscotch,
a trapezoid, a parallelogram, a rectangle, a square,
a house,
a stylized flower,
an origami paper hen,
a little man,
the capital letter H in relief,
the digits of a digital watch,
the alphabet (in capital letters),
a helicopter,
a truck,
a train,
a barge, etc.
$\Box$}\footnote{\/The box $\Box$ marks the end of the text of an example, of an
exercise, of a solution, etc. It allows the reader who does not want to read
this text to find quickly where to resume reading.}\end{exercices}
\begin{figure}[htbp]
 \includevspacePostScript[6.8cm]{completer-les-dessins}{0.95}{0.95}
 \caption{Drawings to complete}
 \label{fig:completer-les-dessins}
\end{figure}
\vfill

\clearpage
\begin{corriges}{\ \\

\vspace{-\baselineskip}
\tiret Once completed, the drawings of figure
\ref{fig:completer-les-dessins} (page \pageref{fig:completer-les-dessins})
rep\-re\-sent a boat and a dog.

\tiret  Examples of drawings
are given on pages \pageref{fig:marelle}, \pageref{fig:quadrilateres} and
\pageref{fig:transports}.
$\Box$\footnote{\ We give the solutions of the exercises a few pages
after their state\-ments.
We strongly advise the reader to read this book in small steps,
taking care to do a few exercises after each chapter.}
}\end{corriges}
\begin{figure}[htbp]
 \includehvspacePostScript[117mm,3.2cm]{marelle}{0.95}{0.95}
 \caption{Hopscotch}
 \label{fig:marelle}
\end{figure}
\vfill
\begin{figure}[htbp]
 \includehvspacePostScript[117mm,8.1cm]{alphabet}{0.95}{0.95}
 \caption{Digital digits and alphabet (in capital letters)}
 \label{fig:alphabet}
\end{figure}
\vfill
\twocolumn
\begin{figure}[h]
 \includevspacePostScript[2.5cm]{quadrilateres}{0.95}{0.95}
 \caption{Quadrilaterals}\label{fig:quadrilateres}
\end{figure}
\begin{figure}[h]
 \includevspacePostScript[5.5cm]{fleur-1}{0.95}{0.95}
 \caption{Stylized flower}
 %\label{fig:fleur-1}
\end{figure}
\vfill
\begin{figure}[h]
 \includevspacePostScript[0cm]{maison}{0.95}{0.95}
 \caption{House}
 %\label{fig:maison}
\end{figure}
\vfil
\begin{figure}[h]
 \includevspacePostScript[2.8cm]{cocotte}{0.95}{0.95}
 \caption{Origami paper hen}
 %\label{fig:cocotte}
\end{figure}
\vfil
\begin{figure}[h]
 \includevspacePostScript[7.6cm]{bonhomme}{0.95}{0.95}
 \caption{Little man}
 %\label{fig:bonhomme}
\end{figure}
\vfil
\begin{figure}[h]
 \includevspacePostScript[4.6cm]{LettreH}{0.95}{0.95}
 \caption{Capital letter H in relief}
 \label{fig:h-relief}
\end{figure}
\onecolumn
\begin{figure}[h]
 \includevspacePostScript[7.4cm]{transports}{0.95}{0.95}
 \caption{Means of transport}\label{fig:transports}
\end{figure}

%%%%%%%%%%%%%%%%%%%%%%%%%%%%%%%%%%%%%%%%%%%%%%%%%%%%%%%%%%%%%%%%%%%%%%%%%%%
\noindent \framebox[\textwidth]{
\begin{minipage}{\myparboxwidth}
\setlength{\parindent}{2em}
\small \sf
\vspace{\abovedisplayskip}
\section*{Addenda}

Before writing a computer program to solve a
problem, one must begin by specifying
precisely the task that the computer will have to carry out. In our
graphic approach to explaining the programming of
com\-put\-ers, this {\em specification} \index{Specification} is a drawing,
which the computer must produce.

Drawing by hand while respecting the constraint of using only
lines of the form
\begin{picture}(12,10)
\put(0,0){\line(1,0){10}}
\end{picture}\/,
\begin{picture}(6,10)
\put(2,0){\line(0,1){10}}
\end{picture}\/,
\begin{picture}(10,10)
\put(0,0){\line(1,1){10}}
\end{picture}\/,
\begin{picture}(10,10)
\put(0,10){\line(1,-1){10}}
\end{picture}\ ,
\begin{picture}(10,10)
\put(10,0){\oval(20,20)[tl]}
\end{picture}\/,
\begin{picture}(10,10)
\put(0,10){\oval(20,20)[br]}
\end{picture}\/,
\begin{picture}(10,10)
\put(0,0){\oval(20,20)[tr]}
\end{picture}\ ,
\begin{picture}(10,10)
\put(10,10){\oval(20,20)[bl]}
\end{picture}\ ,
\begin{picture}(10,10)
\put(0,0){\makebox(10,10)[bl]{$\displaystyle \smile$}}
\end{picture}\ ,
\begin{picture}(10,10)
\put(0,0){\makebox(10,10)[bl]{$\displaystyle \frown$}}
\end{picture}\ ,
\begin{picture}(5,10)
\put(0,0){\makebox(5,10)[bl]{\footnotesize (}}
\end{picture}\/
or
\begin{picture}(5,10)
\put(0,0){\makebox(5,10)[bl]{\footnotesize )}}
\end{picture}\
is difficult for very young children. To ensure a progression from the
simple to the complex, one begins
by having them copy drawings made only of
horizontal and vertical lines on paper with large squares, then one
introduces the oblique segments, the quarter
circles on the diagonals and finally the quarter circles on the
sides of a square of the grid.
These geometric drawings are an excellent preparation for learning to write.
Then the children enjoy inventing their own
drawings. They must then learn to model their ideas
within the rather strict frame set for them by the drawing rules.
This notion of {\em model}, in the sense of a simplified representation of a
complex process or system, is at the basis of the {\em
scientific approach}.

\vspace{\abovedisplayskip}
\end{minipage}
}
%%%%%%%%%%%%%%%%%%%%%%%%%%%%%%%%%%%%%%%%%%%%%%%%%%%%%%%%%%%%%%%%%%%%%%%%%%%%

%%%%%%%%%%%%%%%%%%%%%%%%%%%%%%%%%%%%%%%%%%%%%%%%%%%%%%%%%%%%%%%%%%%%%%%%%%%
\noindent \framebox[\textwidth]{
\begin{minipage}{\myparboxwidth}
\setlength{\parindent}{2em}
\small \sf
\vspace{\abovedisplayskip}
\noindent These geometric drawings, built using
only a few primitive shapes, are
also a first step towards programming, which consists of build\-ing
a sometimes complex edifice from the elementary bricks
that are the instructions of the programming language.

In the next chapter, we no longer make drawings by
hand but draw them on the computer screen by driving
a robot with the computer's keyboard.
This lets us discover how the computer works
in {\em interactive mode} and enter directly into the world of
{\em computer graphics}, \index{Computer graphics} that is, the creation of images on a computer.
By construction, the robot draws while respecting the drawing rules
that have been set. It lets us correct our mistakes easily,
save our drawing on disk, make reductions or
enlargements of it and print as many copies of it as we want.

\vspace{\abovedisplayskip}
\end{minipage}
}
%%%%%%%%%%%%%%%%%%%%%%%%%%%%%%%%%%%%%%%%%%%%%%%%%%%%%%%%%%%%%%%%%%%%%%%%%%%%
\newpage
\thispagestyle{empty}
